Fix \(d\ge4\), \(\eta>0\), and the public experiment \(\mathfrak E\). By \(\text{lem:fixed-experiment-shell-certificate}\), a law in \(\mathcal M_{\mathfrak E}(d,C,D,\eta)\) exists if and only if at least one of the displayed finitely many branches has a feasible triple \((\vartheta,\rho,u)\). Its likelihood-ratio constraints impose \(C_P=C\), its semidefinite and unit-kernel constraints impose \(D_P=D\), and its positive-definiteness constraint imposes \(\Lambda_{\mathrm{ref}}\succ0\). Conversely, every member of the exact shell supplies such a feasible branch by taking \(I\) to be the positive-mass context support, \(\vartheta=\theta\), and \(u\) to be a unit generalized eigenvector for the largest eigenvalue. This proves the fixed-experiment nonemptiness equivalence.

If the boundary shell is nonempty, \(\text{thm:zero-risk-boundary}\) gives \(\pi_P^\star=\pi_{\mathrm{ref}}\) for every member and shows that the constant public learner \(\pi_{\mathrm{ref}}\) has zero regularized regret. Since sample sizes range over \(\mathbb N\),
\[
N^\star_{\mathfrak E}(\epsilon;d,1,1,\eta)=1
\]
for every \(\epsilon>0\).

Now suppose that a nontrivial shell is nonempty. Invoke \(\text{thm:clean-temperature-feature-upper}\). Its learner chooses, inside the data-dependent compact set \(\mathcal Q_n\), a function minimizing the empirical worst Gibbs-weighted squared radius and deploys the associated Gibbs policy. The measurable-maximum construction in that theorem makes the learner a member of \(\mathcal L_n(\mathfrak E)\). Its formula uses only \((\mathcal D_n,\mathfrak E)\), so it uses neither the law-specific \(\rho\) nor the analytical shell labels \(C,D\).

The same theorem gives the deterministic regret cap \(\eta/2\). Therefore, if \(\eta\le2\epsilon\), its risk at \(n=1\) is at most \(\epsilon\), and hence
\[
N^\star_{\mathfrak E}(\epsilon;d,C,D,\eta)=1.
\]
For \(\eta>2\epsilon\), define
\[
q=dD\min\{\eta/\epsilon,\epsilon^{-2}\}.
\]
The explicit inversion in \(\text{thm:clean-temperature-feature-upper}\), including integer rounding and the \(n^{-2}\min\{2,\eta/2\}\) complement contribution, supplies a universal \(K_0\) for which
\[
N^\star_{\mathfrak E}(\epsilon;d,C,D,\eta)
\le
\left\lceil K_0q\left[1+\log(eq)+\log\{1+(1+\eta)q\}\right]\right\rceil
\]
whenever the displayed order is at least two. This is precisely
\[
\widetilde O\!\left(dD\min\{\eta/\epsilon,\epsilon^{-2}\}\right),
\]
where the two displayed logarithms are the only suppressed factors. In particular, the multiplicative constant is universal, the result holds for every \(\eta>0\), and \(C\) does not occur.

We next show why this upper frontier cannot be complemented by an index-only matching lower frontier. Fix any \(1<D\le C<e^\eta\). By \(\text{thm:fixed-design-feasibility-and-frontier-nonuniqueness}\), there is a finite public experiment \(\mathfrak E_0\) with a nonempty exact shell and zero minimax risk for every sample size. Hence
\[
N^\star_{\mathfrak E_0}(\epsilon;d,C,D,\eta)=1
\qquad(\epsilon>0).
\]
The same theorem supplies a finite public experiment \(\mathfrak E_1\) at the identical four indices and constants \(c(C,D,\eta),\epsilon_0(C,D,\eta)>0\) such that
\[
N^\star_{\mathfrak E_1}(\epsilon;d,C,D,\eta)
\ge c(C,D,\eta)\frac{dD\eta}{\epsilon}
\qquad(0<\epsilon\le\epsilon_0(C,D,\eta)).
\]
That construction applies throughout \(1<D\le C<e^\eta\), without the earlier restriction \(2D-1<e^\eta\), and thus includes \(2D-1\ge e^\eta\). Because the leading constant may depend arbitrarily on \((C,D,\eta)\), this is correctly read as a pointwise \(\Omega_{C,D,\eta}(d/\epsilon)\) lower bound. The pair \((\mathfrak E_0,\mathfrak E_1)\) proves that the four indices do not determine a unique fixed-experiment frontier.

It remains to settle the proposed quadratic branch. Fix one finite \(\mathfrak E\), and hold \((d,C,D,\eta)\) fixed. For all sufficiently small \(\epsilon\), one has \(\eta>2\epsilon\) and \(\eta\epsilon<1\), so
\[
q=dD\frac{\eta}{\epsilon}.
\]
The explicit upper bound then implies
\[
N^\star_{\mathfrak E}(\epsilon;d,C,D,\eta)
=O_{d,D,\eta}\!\left(\epsilon^{-1}\log(1/\epsilon)\right),
\]
because each displayed logarithm is \(O_{d,D,\eta}(\log(1/\epsilon))\). Consequently,
\[
0\le
\frac{\epsilon^2}{d}N^\star_{\mathfrak E}(\epsilon;d,C,D,\eta)
=O_{d,D,\eta}\!\left(\epsilon\log(1/\epsilon)\right)
\longrightarrow0.
\]
If constants \(c(C,D,\eta)>0\) and \(\epsilon_1(C,D,\eta)>0\) satisfied
\[
N^\star_{\mathfrak E}(\epsilon;d,C,D,\eta)
\ge c(C,D,\eta)d/\epsilon^2
\]
for every \(0<\epsilon\le\epsilon_1(C,D,\eta)\), then the preceding normalized limit would be at least \(c(C,D,\eta)>0\), a contradiction. Thus the sufficiently-small-accuracy quadratic branch is false for every specified finite experiment.

Finally, let \(L(\epsilon;d,C,D,\eta)\) be any lower bound asserted uniformly over all public experiments having a nonempty exact shell at the four indices. Applying it to \(\mathfrak E_0\) gives
\[
L(\epsilon;d,C,D,\eta)
\le N^\star_{\mathfrak E_0}(\epsilon;d,C,D,\eta)=1.
\]
No nontrivial uniform lower frontier can therefore diverge with \(\epsilon^{-1}\) or \(\epsilon^{-2}\). The clean upper bound depends on \(D\) but not \(C\), proving the asserted achievability irrelevance of pointwise coverage beyond feature coverage. The zero-versus-informative pair simultaneously shows why that irrelevance does not make the public experiment dispensable. This resolves every part of the question.